\section{Question 7}
\begin{enumerate}[label=(\alph*)]
\item
    \begin{enumerate}[label=(\roman*)]
    \item See set sketches below.
    \item 
    Firstly, $A-B$ is bounded since it can be contained in some closed disc. It is also closed. Therefore, it is compact.
    
    The function $f$ is continuous on $A-B$ since it has no singularities within. Therefore, by The Boundedness Theorem, $f$ is bounded on $A-B$.

    Finally, the set $C$ is unbounded, as is $C - \{-1\}$.

    Also, if we write $z+i = x+iy$,
    \[
    \implies \exp\left(\frac{i}{z+i}\right) = \exp\left(\frac{y}{x^2+y^2}\right)\exp\left(i\frac{x}{x^2+y^2}\right)
    \]
    it is clear that as $y\rightarrow\infty$, $f\rightarrow\infty$. Hence, $f$ is unbounded on $C-\{-1\}$.

    \textcolor{purple}{I think this is valid reasoning, since as $y\rightarrow\infty$ the imaginary part $\rightarrow e^{i/x}$ and the real part $\rightarrow \infty$, so it's a circle tending towards infinite radius which is clearly unbounded. However, I'm not 100\% sure. I'm certain both of you did it differently, we can consider replacing what I did with one of your solutions Donach and John?\\ -- Thomas}

    \textcolor{blue}{This is the only question on the paper I didn't do. I'll have a think about it when I've time.\\ -- John}
    
    \textcolor{teal}{The edit I'd make to your solution is to add modulus signs, so as $ y \to \infty, |f(z)| \to \infty $. 
    \vs \\
    I'd also wonder if it's worth including a sequence of points, as in the solution to Spec 1 Q7(a)(ii)? \\
    Eg, for $n \in \N, \; (n+2)i \in C$, and  $|f\big((n+2)i\big)| = \left| \exp\dfrac{n+3}{(n+3)^2} \right| \left| \exp i(0) \right| = \exp(n+3) \to \infty $ as $ n \to \infty $, 
    \\ -- Donach
    \vs \\
    Hang on, that doesn't work. As $y \to \infty $, the real part goes to $ e^0 = 1 $.
    \vs \\
    I'll add a solution based on my answer.
    }
    \vs \\
    Finally, $f$ is unbounded on the set $C - \{-1\}$.
    \vs \\
    As $f$ has a pole at -1, by HB B4 3.3 p58, $ f(z) \to \infty $, as $ z \to -1 $, e.g. along the points $ -1 + \dfrac{1}{2n}, \; n=1,2,\ldots $, which are all in $ C - \{  -1 \} $.
    \vs \\
    {
    \color{teal}
    Oops, I think that's wrong as well. It should be the CW theorem as it's an essential singularity. I'm going to push this and then do another version.
    }
    \end{enumerate}
    \leavevmode \vvs \\

    \item 
    \begin{enumerate}
        \item 
        Let,
        \[
        z = x + iy
        \]
        \[
        \implies \overline{z} = x - iy.
        \]
        Then,
        \[
        h(x+iy) = (x-iy)^2 + (x+iy)(x-iy) - y
        \]
        \[
        = x^2 + i^2y^2 - 2ixy +x^2 +y^2 - y
        \]
        \[
        = 2x^2 - y - 2ixy,
        \]
        as required.
        \item 
        We let $h(x+iy)=u(x,y)+iv(x,y)$ be defined on a region $\C$ containing $a+ib$. Then,
        \[
        u(x,y)=2x^2 - y,
        \]
        and
        \[
        v(x,y)=-2xy
        \]
        and hence their partials are,
        \[
        u_x = 4x,\quad u_y = -1,
        \]
        and
        \[
        v_x = -2y,\quad v_y = -2x.
        \]
        We now check the point $(a,b)$ via the Cauchy-Riemann equations,
        \[
        u_x(a,b) = 4a = v_y(a,b) = -2a
        \]
        \[
        \implies 4a = -2a
        \]
        \[
        \implies a = 0,
        \]
        and
        \[
        v_x(a,b)=-2b = -u_y(a,b) = -1
        \]
        \[
        \implies b = \frac{1}{2}.
        \]
        Therefore, the CR equations are satisfied at $(0, \frac{1}{2})$, and since the partial derivatives exist at $(x,y)$ for each $x+iy \in \C$ and are continuous at $(0, \frac{1}{2})$, $h$ is differentiable at this point and has a derivative here of,
        \[
        h'\left(0+\frac{1}{2}i\right) = - i,
        \]
        by the Cauchy-Riemann Converse Theorem. However, the Cauchy-Riemann equations fail at points other than $\left(0, \frac{1}{2}\right)$, so $f$ is not differentiable on $ \C - \{\left(0, \frac{1}{2}\right)\}$ by the Cauchy-Riemann Theorem.
        \end{enumerate}
        {
        \color{teal}
        I also added in a sentence explicitly referencing the CR Theorem. I think there's a mark for that.
        }
\end{enumerate}

\begin{figure}[p]
\centering
    % TIKZ PICTURE FOR Q7-ai
    \input{tikz/Q7/Q7-ai-set-sketches}
\end{figure}
